%\begin{table}[htb]
%\centering
%\begin{tabular}{| c | c | c | c | c |}
%\hline
%& \multicolumn{2}{c|}{\bf electron-like} & \multicolumn{2}{c|}{\bf muon-like} \\
%\hline
%\bf Source of uncertainty			& \bf neutrino  	& \bf antineutrino  	& \bf neutrino  	& \bf antineutrino	\\
%\hline
%  Super-K detector                   		& 2.63\% 		& 3.66\%  			& 3.87\%  		& 3.25\%      	  	\\
%  FSI+SI+PN                 			& 2.54\% 		& 2.44\% 			& 1.44\% 		& 2.06\%                   	\\
%  Flux                           			& 3.09\%  		& 3.26\%  			& 2.98\%  		& 3.08\%              	\\
%  Cross section 		               		& 4.81\%    	& 5.19\%    		& 3.70\%    	& 3.74\%                   	\\
%  Flux+Cross section                      	& 3.85\%    	& 4.17\%    		& 2.35\%    	& 2.40\%                   	\\
%  Oscillations                   			& 4.24\%   	& 4.02\%   		& 0.02\%   	& 0.03\%                   	\\
%  \hline
%  Total                            			& 6.75\%  		& 7.41\%  			& 4.73\%  		& 4.53\%                   	\\
%\hline
%\end{tabular}


%%% fhc mu-like

\begin{table}[htb]
\centering
\begin{tabular}{| c | c | c |}
\hline
  {\bf Source of uncertainty}    & \multicolumn{1}{c|}{{\bf ${\delta}N_{SK}$/$N_{SK}$}}\\
                                 & \multicolumn{1}{c|}{(Asimov A)}\\

\hline
\hline
SKDet+FSI+SI                   & 4.15\%                        \\
\hline
\hline
  Flux                           & 3.59\%                        \\
\hline
\hline
  2p-2h (corr)                     & 3.54\%                        \\
  2p-2h-bar (corr)                  & 0.20\%                        \\
  NC other (uncorr)              & 0.77\%                        \\
  NC 1gamma (uncorr)             & 0.00\%                        \\
  XSec nue/numu (uncorr)                 & 0.01\%                        \\
\hline
  XSec Tot (corr)                  & 4.02\%                        \\
  XSec Tot                       & 4.10\%                        \\
\hline
\hline
  Flux+XSec (ND280 constrained)                      & 2.82\%                        \\
  Flux+XSec  (All)                    & 2.92\%                        \\
\hline
  Flux+XSec+SKDet+FSI+SI         & 5.09\%                        \\
\hline
\hline
  Oscillations                   & 0.03\%                        \\
\hline
\hline
  All                            & 5.13\%                        \\
    All (pre-fit) 		& 11.9\%				\\
\hline

\end{tabular}
\caption{
 Effect of 1$\sigma$ variations on the total number of events 1-ring muon-like events for neutrino mode
 given by all the systematic uncertainties, obtained by performing 10k toys MC.
The RMS / Mean of the distribution is assumed to be the $1\sigma$ uncertainty.
The systematic parameter values measured with the ND280 data 
shown in 
\cref{tab:banff_uncertainties_flux_fhc},
~\ref{tab:banff_uncertainties_flux_rhc}
~\ref{tab:banff_uncertainties_xsec}
and
~\ref{tab:banff_uncertainties_skdetfsi}
and the oscillation parameters of the Asimov data set A shown in \cref{tab:osc_params_joint_asimov}
are used.
Since the parameters $\delta_{CP}$, $\theta_{23}$ and $\Delta m^2_{32}$ ($\Delta m^2_{13}$) and mass hierarchy
will be measured,
only the uncertainty on the parameters 
$\sin^2 \theta_{13}$,
$\sin^2 \theta_{12}$
and 
$\Delta m^2_{21}$ 
are taken into account.
Pre-BANFF fit systematic uncertainties are used were stated.
}                                                                                                                                                                                          
\label{tab:nsk_systematic_table_summary_joint_asimov_A_fhc_mulike}
\end{table}




%%%%% fhc e-like


\begin{table}[htb]
\centering
\begin{tabular}{| c | c | c |}
\hline
  {\bf Source of uncertainty}    & \multicolumn{1}{c|}{{\bf ${\delta}N_{SK}$/$N_{SK}$}}\\
                                 & \multicolumn{1}{c|}{(Asimov A)}\\

\hline
\hline  SKDet+FSI+SI                   & 3.50\%                        \\
\hline
\hline
  Flux                           & 3.67\%                        \\
\hline
\hline
  2p-2h (corr)                     & 3.85\%                        \\
  2p-2h bar (corr)                  & 0.05\%                        \\
  NC other (uncorr)              & 0.16\%                        \\
  NC 1gamma (uncorr)             & 1.47\%                        \\
  XSec nue/numu (uncorr)                 & 2.60\%                        \\
\hline
  XSec Tot (corr)                  & 4.25\%                        \\
  XSec Tot                       & 5.20\%                        \\
\hline
\hline
  Flux+XSec (ND280 constrained)                      & 2.84\%                        \\
  Flux+XSec  (All)                     & 4.13\%                        \\
\hline
  Flux+XSec+SKDet+FSI+SI         & 5.53\%                        \\
\hline
\hline
  Oscillations                   & 4.16\%                        \\
\hline
\hline
  All                            & 6.80\%                        \\
  All (pre-fit) 		& 12.6\%				\\
\hline
\hline

\end{tabular}
\caption{
 Effect of 1$\sigma$ variations on the total number of events 1-ring electron-like events for neutrino mode
 given by all the systematic uncertainties, obtained by performing 10k toys MC.
The RMS / Mean of the distribution is assumed to be the $1\sigma$ uncertainty.
The systematic parameter values measured with the ND280 data 
shown in 
\cref{tab:banff_uncertainties_flux_fhc},
~\ref{tab:banff_uncertainties_flux_rhc}
~\ref{tab:banff_uncertainties_xsec}
and
~\ref{tab:banff_uncertainties_skdetfsi}
and the oscillation parameters of the Asimov data set A shown in \cref{tab:osc_params_joint_asimov}
are used.
Since the parameters $\delta_{CP}$, $\theta_{23}$ and $\Delta m^2_{32}$ ($\Delta m^2_{13}$) and mass hierarchy
will be measured,
only the uncertainty on the parameters 
$\sin^2 \theta_{13}$,
$\sin^2 \theta_{12}$
and 
$\Delta m^2_{21}$ 
are taken into account.
Pre-BANFF fit systematic uncertainties are used were stated.
}                                                                                                                                                                                          
\label{tab:nsk_systematic_table_summary_joint_asimov_A_fhc_elike}
\end{table}



%%%%%% rhc mu-like


\begin{table}[htb]
\centering
\begin{tabular}{| c | c | c |}
\hline
  {\bf Source of uncertainty}    & \multicolumn{1}{c|}{{\bf ${\delta}N_{SK}$/$N_{SK}$}}\\
                                 & \multicolumn{1}{c|}{(Asimov A)}\\

\hline
\hline  SKDet+FSI+SI                   & 3.94\%                        \\
\hline
  Flux                           & 3.66\%                        \\
\hline
\hline
  2p-2h (corr)                     & 2.96\%                        \\
  2p-2h bar (corr)                  & 1.82\%                        \\
  NC other (uncorr)              & 0.75\%                        \\
  NC 1gamma (uncorr)             & 0.00\%                        \\
  XSec nue/numu (uncorr)                 & 0.00\%                        \\
\hline
  XSec Tot (corr)                  & 4.08\%                        \\
  XSec Tot                       & 4.14\%                        \\
\hline
\hline
  Flux+XSec (ND280 constrained)                      & 3.27\%                        \\
  Flux+XSec  (All)                    & 3.36\%                        \\
\hline
  Flux+XSec+SKDet+FSI+SI         & 5.18\%                        \\
\hline
\hline
  Oscillations                   & 0.03\%                        \\
\hline
\hline
  All                            & 5.12\%                        \\
      All (pre-fit) 		& 12.7\%				\\
\hline
\hline
\end{tabular} 
\caption{
 Effect of 1$\sigma$ variations on the total number of events 1-ring muon-like events for antineutrino mode
 given by all the systematic uncertainties, obtained by performing 10k toys MC.
The RMS / Mean of the distribution is assumed to be the $1\sigma$ uncertainty.
The systematic parameter values measured with the ND280 data 
shown in 
\cref{tab:banff_uncertainties_flux_fhc},
~\ref{tab:banff_uncertainties_flux_rhc}
~\ref{tab:banff_uncertainties_xsec}
and
~\ref{tab:banff_uncertainties_skdetfsi}
and the oscillation parameters of the Asimov data set A shown in \cref{tab:osc_params_joint_asimov}
are used.
Since the parameters $\delta_{CP}$, $\theta_{23}$ and $\Delta m^2_{32}$ ($\Delta m^2_{13}$) and mass hierarchy
will be measured,
only the uncertainty on the parameters 
$\sin^2 \theta_{13}$,
$\sin^2 \theta_{12}$
and 
$\Delta m^2_{21}$ 
are taken into account.
Pre-BANFF fit systematic uncertainties are used were stated.
}                                                                                                                                                                                          
\label{tab:nsk_systematic_table_summary_joint_asimov_A_rhc_mulike}
\end{table}



%%%%%% rhc e-like

\begin{table}[htb]
\centering
\begin{tabular}{| c | c | c |}
\hline
  {\bf Source of uncertainty}    & \multicolumn{1}{c|}{{\bf ${\delta}N_{SK}$/$N_{SK}$}}\\
                                 & \multicolumn{1}{c|}{(Asimov A)}\\

\hline
\hline  SKDet+FSI+SI                   & 3.97\%                        \\
\hline
  Flux                           & 3.78\%                        \\
\hline
\hline
  2p-2h (corr)                     & 3.00\%                        \\
  2p-2h bar (corr)                  & 2.36\%                        \\
  NC other (uncorr)              & 0.33\%                        \\
  NC 1gamma (uncorr)             & 2.93\%                        \\
  XSec nue/numu (uncorr)                 & 1.51\%                        \\
\hline
  XSec Tot (corr)                  & 4.30\%                        \\
  XSec Tot                       & 5.43\%                        \\
\hline
\hline
  Flux+XSec (ND280 constrained)                      & 3.19\%                        \\
  Flux+XSec                      & 4.61\%                        \\
\hline
  Flux+XSec+SKDet+FSI+SI         & 6.21\%                        \\
\hline
\hline
  Oscillations                   & 4.00\%                        \\
\hline
\hline
  All                            & 7.41\%                        \\
      All (pre-fit) 		& 14.3\%				\\
\hline
\hline
\end{tabular}
\caption{
 Effect of 1$\sigma$ variations on the total number of events 1-ring electron-like events for antineutrino mode
 given by all the systematic uncertainties, obtained by performing 10k toys MC.
The RMS / Mean of the distribution is assumed to be the $1\sigma$ uncertainty.
The systematic parameter values measured with the ND280 data 
shown in 
\cref{tab:banff_uncertainties_flux_fhc},
~\ref{tab:banff_uncertainties_flux_rhc}
~\ref{tab:banff_uncertainties_xsec}
and
~\ref{tab:banff_uncertainties_skdetfsi}
and the oscillation parameters of the Asimov data set A shown in \cref{tab:osc_params_joint_asimov}
are used.
Since the parameters $\delta_{CP}$, $\theta_{23}$ and $\Delta m^2_{32}$ ($\Delta m^2_{13}$) and mass hierarchy
will be measured,
only the uncertainty on the parameters 
$\sin^2 \theta_{13}$,
$\sin^2 \theta_{12}$
and 
$\Delta m^2_{21}$ 
are taken into account.
Pre-BANFF fit systematic uncertainties are used were stated.
}                                                                                                                                                                                          
\label{tab:nsk_systematic_table_summary_joint_asimov_A_rhc_elike}
\end{table}


















\begin{table}
\center
\begin{tabular}{| c || c | c || c | c |}
\hline
Source of uncertainty 
&$\mu$-like ${\delta} \Big( \frac{\# {\nu} \text{-mode}}{\# {\bar{\nu}} \text{-mode} } \Big) $ / $\Big \langle \frac{\# {\nu} \text{-mode}}{\# {\bar{\nu}} \text{-mode} } \Big \rangle$ 
&$e$-like ${\delta} \Big( \frac{\# {\nu} \text{-mode}}{\# {\bar{\nu}} \text{-mode} } \Big) $ / $\Big \langle \frac{\# {\nu} \text{-mode}}{\# {\bar{\nu}} \text{-mode} } \Big \rangle$ \\


\hline
SKDet & 0.7\% & 1.6\% \\ \hline
FSI+SI & 2.6\% & 3.5\% \\ \hline
SKDet+FSI+SI & 2.6\% & 3.8\% \\ \hline
\hline
Flux & 1.7\% & 1.7\% \\ \hline
XSec NC other (uncorr) & 0.0\% & 0.2\% \\ \hline
XSec NC $1\gamma$ (uncorr)  & 0.0\%  & 1.5\% \\ \hline
XSec $\nu_e$ / $\nu_{\mu}$ (uncorr)  & 0.0\%  & 3.1\% \\ \hline
XSec Tot  (All) & 1.7\%  & 4.2\% \\ \hline
 \hline
Flux+XSec  & 1.8\%  & 4.1\% \\ \hline
 \hline
Flux+XSec+SDet+FSI+SI  & 3.2\%  & 5.7\% \\ \hline
 \hline
Oscillations  & 0.0\%  & 0.8\% \\ \hline
 \hline
All  & 3.2\%  & 5.8\% \\ \hline
\end{tabular}
\caption{
 Effect of 1$\sigma$ variations of the several source of systematic uncertainties 
 on the ratio between the
 total number of events in neutrino and antineutrino modes
 for 1-ring electron- and muon-like samples
 obtained by performing 10k toys MC.
 The systematic parameter values measured with the ND280 data 
shown in 
\cref{tab:banff_uncertainties_flux_fhc},
~\ref{tab:banff_uncertainties_flux_rhc}
~\ref{tab:banff_uncertainties_xsec}
and
~\ref{tab:banff_uncertainties_skdetfsi}
and the oscillation parameters of the Asimov data set A shown in \cref{tab:osc_params_joint_asimov}
are used.
Since the parameters $\delta_{CP}$, $\theta_{23}$ and $\Delta m^2_{32}$ ($\Delta m^2_{13}$) and mass hierarchy
will be measured,
only the uncertainty on the parameters 
$\sin^2 \theta_{13}$,
$\sin^2 \theta_{12}$
and 
$\Delta m^2_{21}$ 
are taken into account.
}             
\label{tab:ratio_systematic_table_summary_asimov_A}                                                                                                                                                                             
\end{table}
